\section{Bulgular}\label{sec:bulgular}

Ana sonuçlar, Aşama~2 ve 3'ün hiçbir geliştirme koşusunda kullanılmamış 54 problemden (holdout2),
10 tohum ve 10~sn bütçeyle elde edildi. Farklar gap puanıdır; negatif değer ilk stratejinin lehinedir.
Bölüm önce ön-kayıtlı doğrulayıcı testi, ardından keşifsel katman, mekanizma, duyarlılık ve
sağlamlık analizlerini, sonra yayımlanmış bir öğrenilmiş modelle, eşit vuruş sayısıyla ve LKH-3 ile
yapılan hakem odaklı blokları verir. Son alt bölüm, Aşama~3 hipotezlerini doğuran Aşama~2 sonuçlarını
özetler: Aşama~2'de eğitimsiz kNN sırasının GNN'yi geçmesi, eğitimsiz bir vekilin yeterli olabileceği
fikrini doğurdu; Aşama~3 bu fikri yeni bir test setinde ön-kayıtlı olarak sınadı.
% kaynak: E586-E588

\subsection{Doğrulayıcı test}\label{sec:bulgular-dogrulayici}

Sabit sırayla sınanan dört ön-kayıtlı hipotezin dördü de desteklendi
(Tablo~\ref{tab:v3-confirm}). Koşu içinde ağırlıklarını güncelleyen çevrimiçi model, aynı ağırlıklarla
başlayıp hiç güncellenmeyen sabit doğrusal modelden $0{,}291$ puan kötüdür ve 54 problemin hiçbirinde
onu geçemez (V1, $p=1{,}6\times10^{-10}$).
Eğitimsiz fazla-uzunluk vekili, rastgele vuruşa göre ortalama gap'i $0{,}172$ puan düşürür
(V2, $p=3{,}7\times10^{-7}$).
Vekil, LKH-3 turlarıyla eğitilmiş GNN'ye (V3) ve kazanç hedefiyle eğitilmiş MLP'ye (V4) karşı
$0{,}05$ puanlık sınırla kötü-değil testlerini geçer; gözlenen farklar vekil lehinedir ($-0{,}119$
ve $-0{,}083$ puan; $p=7{,}9\times10^{-9}$ ve $1{,}8\times10^{-7}$).
Problem başı farklar sağa çarpık olduğundan Hodges--Lehmann kestirimleri ortalamalardan küçüktür
(V1 $+0{,}242$ [$+0{,}167$; $+0{,}337$], V2 $-0{,}101$ [$-0{,}166$; $-0{,}062$], V3 $-0{,}069$
[$-0{,}114$; $-0{,}039$], V4 $-0{,}057$ [$-0{,}082$; $-0{,}032$]); kararlar değişmez, V3
ve V4'te üst sınırlar ($-0{,}039$ ve $-0{,}032$) sıfırın altındadır.
Aşama~1'de kullanılmış üç TSPLIB problemi çıkarıldığında da (51 problem; keşifsel duyarlılık
denetimi) dört hipotezin dördü reddedilir (V1 $+0{,}295$, V2 $-0{,}166$, V3 $-0{,}114$, V4 $-0{,}084$
puan; en büyük $p=7{,}5\times10^{-7}$).
Ön-kayıtlı iddia bu dört sonuçtan ibarettir. Aynı tablodaki Holm düzeltmeli ikincil çiftler
keşifseldir: vekil GNN'den ($p_{\mathrm{Holm}}=1{,}0\times10^{-4}$) ve kNN sırasından ($-0{,}093$ puan,
$p_{\mathrm{Holm}}=3{,}4\times10^{-5}$) iyidir; GNN ile kNN ayrışmaz ($p_{\mathrm{Holm}}=1{,}0$).
% kaynak: E589, E599-E602, E604, E605, E607, E960-E963, E1100-E1103

\begin{table}[!htbp]
\centering\small\setlength{\tabcolsep}{3pt}
\caption{Doğrulayıcı test (PROTOCOL\_v2, holdout2: 54 problem, 10 tohum, 10~sn). $\Delta$: A$-$B gap farkı [puan], negatif A lehine; kötü-değil sınırı 0{,}05. Renk: anlamlı fark. Ö/E/G: A'nın önde, eşit ve geride olduğu problem sayısı.}
\label{tab:v3-confirm}
\begin{tabularx}{\textwidth}{@{}l>{\raggedright\arraybackslash}Xlccccl@{}}
\toprule
ID & Karşılaştırma (A$-$B) & Test & $\Delta$ & \%95 GA & Ö/E/G & $p$ & Karar \\
\midrule
V1 & Çevrimiçi $-$ Sabit doğrusal & üstünlük & \kotu{+0{,}291} & [+0{,}221; +0{,}367] & 0/0/54 & $1{,}6\times10^{-10}$ & $H_0$ ret \\
V2 & Fazla-uzunluk $-$ Rastgele & üstünlük & \iyi{$-$0{,}172} & [$-$0{,}255; $-$0{,}098] & 44/0/10 & $3{,}7\times10^{-7}$ & $H_0$ ret \\
V3 & Fazla-uzunluk $-$ GNN & kötü-değil & \iyi{$-$0{,}119} & [$-$0{,}195; $-$0{,}051] & 41/0/13 & $7{,}9\times10^{-9}$ & $H_0$ ret \\
V4 & Fazla-uzunluk $-$ Kazanç-MLP & kötü-değil & \iyi{$-$0{,}083} & [$-$0{,}146; $-$0{,}026] & 42/0/12 & $1{,}8\times10^{-7}$ & $H_0$ ret \\
\midrule
\multicolumn{8}{l}{\textit{İkincil, keşifsel (protokolde adı geçen 5 çift; $p$ Holm düzeltmeli)}} \\
-- & kNN $-$ Rastgele & üstünlük & $-$0{,}079 & [$-$0{,}149; $-$0{,}018] & 31/0/23 & 0{,}169 & anlamsız \\
-- & Fazla-uzunluk $-$ kNN & üstünlük & \iyi{$-$0{,}093} & [$-$0{,}144; $-$0{,}040] & 43/0/11 & $3{,}4\times10^{-5}$ & anlamlı \\
-- & GNN $-$ kNN & üstünlük & +0{,}026 & [$-$0{,}027; +0{,}089] & 30/0/24 & 1{,}000 & anlamsız \\
-- & Kazanç-MLP $-$ Üyelik-MLP & üstünlük & $-$0{,}006 & [$-$0{,}052; +0{,}038] & 25/0/29 & 1{,}000 & anlamsız \\
-- & Fazla-uzunluk $-$ GNN & üstünlük & \iyi{$-$0{,}119} & [$-$0{,}195; $-$0{,}051] & 41/0/13 & $1{,}0\times10^{-4}$ & anlamlı \\
\bottomrule
\end{tabularx}
\end{table}


Şekil~\ref{fig:manset}'teki problem düzeyi dağılımı çarpıktır: vekil$-$GNN farkının medyanı
$-0{,}060$ puandır. Ham kazanma oranı 41/54'tür; ölçüm gürültüsünün (0,016 puan) altındaki farklar
beraberlik sayıldığında oran \%70'e iner. Şekilde işaretli örnek problemler, GNN'nin kendi
çeyrekler arası aralığında kaldığı problemler arasından katman başına en büyük avantajla seçildi.
% kaynak: E789, V013, V014

\subsection{Strateji sıralaması ve problem türleri}\label{sec:bulgular-katman}

Tablo~\ref{tab:v3-strategies} sekiz stratejiyi dört problem türünde karşılaştırır. Fazla-uzunluk,
tüm sette en düşük ortalama gap'i verir (\%0,922; rastgele \%1,094, GNN \%1,041) ve dört katmanın
hepsinde sıralamanın başındadır.
Katman sonuçları betimseldir; tablodaki renkler katman içi Holm düzeltmesine dayanır. Hedeflemenin getirisi problem
yapısına bağlıdır: rastgeleye göre kazanç gerçek problemlerde $0{,}323$, kümeli problemlerde
$0{,}302$ puandır; karışık katmanda $0{,}052$ puana iner, uniform katmanda $0{,}030$ puanda kalır ve
anlamlı değildir ($p=0{,}12$).
kNN sırası bu ayrımda kararsızdır. Kümeli problemlerde rastgeleden $0{,}271$ puan iyi, karışık
problemlerde $0{,}039$ puan kötüdür (düzeltmesiz, betimsel) ve tüm sette rastgeleden anlamlı ayrışmaz
($p_{\mathrm{Holm}}=0{,}17$).
Öğrenilmiş üç skor (GNN, kazanç-MLP, üyelik-MLP) tüm sette rastgeleden iyidir ama vekilin gerisinde
kalır; çevrimiçi model her katmanda en kötüdür.
% kaynak: E603, E608, E610-E613, E615, E618, E626, E633, E634, E641, E642

Yapıya bağımlılık, katman etiketinden bağımsız bir ölçüyle de görülür (keşifsel). Her problem için
en yakın komşu uzaklığının değişim katsayısını (CV) bir kümelenme göstergesi olarak kullandık. Bu
değer uniform problemlerde $0{,}51$ ile $0{,}55$ arasındadır, kümeli ve ulusal problemlerde daha
büyüktür.
Vekilin rastgeleye göre kazancı CV arttıkça büyür (Spearman $\rho=-0{,}56$, $p=9{,}8\times10^{-6}$;
Şekil~\ref{fig:v4-structure}). GNN'ye karşı farkta ilişki aynı yöndedir ama daha zayıftır
($\rho=-0{,}38$, $p=0{,}005$); problem boyutuyla ilişki yoktur ($\rho=-0{,}10$, $p=0{,}46$).
Ölçü, doğrulayıcı koşunun mevcut verisi üzerinde hesaplandığı için betimseldir. Uygulayıcı için
sonucu şudur: hedeflemeden en çok, kenar uzunluklarının düzensiz olduğu kümeli ve gerçek problemlerde
yararlanılır; bu ölçü koşudan önce, koordinatlardan ucuzca hesaplanabilir.
% kaynak: E1038-E1041, E1043, E1045

\begin{figure}[!htbp]
\centering
% Problem yapisi ve hedefleme kazanci -- OTOMATIK URETIM: evidence/scripts/build_v4.py
% Veri: figures/data/v4_structure.dat (v4_robustness_analysis.json T_structure.rows)
\begin{tikzpicture}
\begin{axis}[width=\textwidth, height=5.6cm, xmode=log, log basis x=10, xmin=0.25, xmax=2.5, ymax=0.9,
  xtick={0.3,0.5,0.75,1,1.5,2}, xticklabels={0.3,0.5,0.75,1,1.5,2},
  xlabel={Clustering index: CV of nearest-neighbour distance (log scale)},
  ylabel={Excess $-$ random [gap points]}, tick label style={font=\scriptsize}, label style={font=\small},
  ymajorgrids, grid style={black!10},
  legend style={font=\scriptsize, draw=black!40, fill=white, at={(0.01,0.03)}, anchor=south west, legend columns=4}]
\addplot[black!50, thin, forget plot] coordinates {(0.25,0) (2.5,0)};
\addplot[only marks, mark=x, mark size=1.8pt, black!45] table[x=cv, y=d_rand, restrict expr to domain={\thisrow{guni}}{1:1}] {figures/data/v4_structure.dat};
\addlegendentry{Uniform}
\addplot[only marks, mark=square*, mark size=1.3pt, black!70] table[x=cv, y=d_rand, restrict expr to domain={\thisrow{gmix}}{1:1}] {figures/data/v4_structure.dat};
\addlegendentry{Mixed}
\addplot[only marks, mark=*, mark size=1.5pt, iyiRenk] table[x=cv, y=d_rand, restrict expr to domain={\thisrow{gclu}}{1:1}] {figures/data/v4_structure.dat};
\addlegendentry{Clustered}
\addplot[only marks, mark=triangle*, mark size=1.9pt, kotuRenk] table[x=cv, y=d_rand, restrict expr to domain={\thisrow{greal}}{1:1}] {figures/data/v4_structure.dat};
\addlegendentry{Real}
\node[anchor=north west, font=\scriptsize, fill=white, inner sep=2pt] at (rel axis cs:0.01,0.97) {Spearman $\rho=-0.56$, $p=9.8\times10^{-6}$};
\end{axis}
\end{tikzpicture}

\caption{Problem yapısı ve hedefleme kazancı (holdout2, 10~sn, keşifsel): yatay eksen en yakın komşu uzaklığının değişim katsayısı, dikey eksen fazla-uzunluk vekilinin rastgeleye göre problem başı gap farkı (negatif = vekil iyi).}
\label{fig:v4-structure}
\end{figure}

% OTOMATIK URETIM -- evidence/scripts/build_v3.py (elle duzenleme)
\begin{table}[!htbp]
\centering\small\setlength{\tabcolsep}{3pt}
\caption{Holdout2'de sekiz strateji, katman bazında: ortalama gap ve rastgeleye eşlenik fark $\Delta$. Renk: gap sütununda katmanın en iyi/en kötüsü, $\Delta$'da katman içi Holm düzeltmeli $p<0{,}05$; $p$ sütunu düzeltmesizdir.}
\label{tab:v3-strategies}
\begin{tabular*}{\textwidth}{@{\extracolsep{\fill}}lccccc@{}}
\toprule
Strateji & Gap [\%] & $\Delta$ [puan] & \%95 GA & K/B/K & $p$ \\
\midrule
\textit{Tümü (54 problem)}: Rastgele & 1{,}094 & \multicolumn{4}{c}{referans} \\ % E608
\quad kNN (0,04/0,75) & 1{,}015 & $-$0{,}079 & [$-$0{,}149; $-$0{,}018] & 31/0/23 & 0{,}056 \\ % E609
\quad Fazla-uzunluk (0,06/0,5) & \iyi{0{,}922} & \iyi{$-$0{,}172} & [$-$0{,}255; $-$0{,}098] & 44/0/10 & $3{,}7\times10^{-7}$ \\ % E610
\quad GNN (0,04/0,5) & 1{,}041 & \iyi{$-$0{,}053} & [$-$0{,}080; $-$0{,}026] & 37/0/17 & $6{,}8\times10^{-4}$ \\ % E611
\quad Kazanç-MLP & 1{,}005 & \iyi{$-$0{,}089} & [$-$0{,}152; $-$0{,}037] & 36/0/18 & 0{,}003 \\ % E612
\quad Üyelik-MLP & 1{,}011 & \iyi{$-$0{,}083} & [$-$0{,}127; $-$0{,}044] & 35/0/19 & $6{,}6\times10^{-4}$ \\ % E613
\quad Sabit doğrusal & 1{,}012 & \iyi{$-$0{,}082} & [$-$0{,}157; $-$0{,}026] & 34/0/20 & 0{,}020 \\ % E614
\quad Çevrimiçi (ılık) & \kotu{1{,}303} & \kotu{+0{,}209} & [+0{,}150; +0{,}271] & 4/0/50 & $1{,}4\times10^{-8}$ \\ % E615
\addlinespace
\textit{Gerçek (12 problem)}: Rastgele & 1{,}829 & \multicolumn{4}{c}{referans} \\ % E616
\quad kNN (0,04/0,75) & 1{,}724 & $-$0{,}104 & [$-$0{,}280; +0{,}048] & 9/0/3 & 0{,}151 \\ % E617
\quad Fazla-uzunluk (0,06/0,5) & \iyi{1{,}505} & \iyi{$-$0{,}323} & [$-$0{,}506; $-$0{,}159] & 10/0/2 & 0{,}002 \\ % E618
\quad GNN (0,04/0,5) & 1{,}742 & $-$0{,}086 & [$-$0{,}170; $-$0{,}011] & 8/0/4 & 0{,}077 \\ % E619
\quad Kazanç-MLP & 1{,}715 & $-$0{,}114 & [$-$0{,}240; $-$0{,}006] & 9/0/3 & 0{,}092 \\ % E620
\quad Üyelik-MLP & 1{,}704 & $-$0{,}125 & [$-$0{,}213; $-$0{,}034] & 9/0/3 & 0{,}042 \\ % E621
\quad Sabit doğrusal & 1{,}754 & $-$0{,}075 & [$-$0{,}170; +0{,}010] & 8/0/4 & 0{,}233 \\ % E622
\quad Çevrimiçi (ılık) & \kotu{2{,}157} & \kotu{+0{,}328} & [+0{,}208; +0{,}460] & 1/0/11 & $9{,}8\times10^{-4}$ \\ % E623
\addlinespace
\textit{Uniform (14 problem)}: Rastgele & 0{,}606 & \multicolumn{4}{c}{referans} \\ % E624
\quad kNN (0,04/0,75) & 0{,}625 & +0{,}019 & [$-$0{,}005; +0{,}044] & 6/0/8 & 0{,}241 \\ % E625
\quad Fazla-uzunluk (0,06/0,5) & \iyi{0{,}576} & $-$0{,}030 & [$-$0{,}057; $-$0{,}003] & 10/0/4 & 0{,}119 \\ % E626
\quad GNN (0,04/0,5) & 0{,}590 & $-$0{,}016 & [$-$0{,}050; +0{,}019] & 8/0/6 & 0{,}326 \\ % E627
\quad Kazanç-MLP & 0{,}607 & +0{,}001 & [$-$0{,}028; +0{,}027] & 7/0/7 & 0{,}715 \\ % E628
\quad Üyelik-MLP & 0{,}593 & $-$0{,}013 & [$-$0{,}041; +0{,}016] & 8/0/6 & 0{,}502 \\ % E629
\quad Sabit doğrusal & 0{,}592 & $-$0{,}014 & [$-$0{,}045; +0{,}018] & 9/0/5 & 0{,}463 \\ % E630
\quad Çevrimiçi (ılık) & \kotu{0{,}712} & \kotu{+0{,}106} & [+0{,}082; +0{,}131] & 0/0/14 & $1{,}2\times10^{-4}$ \\ % E631
\addlinespace
\textit{Kümeli (14 problem)}: Rastgele & 1{,}400 & \multicolumn{4}{c}{referans} \\ % E632
\quad kNN (0,04/0,75) & 1{,}129 & \iyi{$-$0{,}271} & [$-$0{,}452; $-$0{,}131] & 13/0/1 & $8{,}5\times10^{-4}$ \\ % E633
\quad Fazla-uzunluk (0,06/0,5) & \iyi{1{,}098} & \iyi{$-$0{,}302} & [$-$0{,}533; $-$0{,}097] & 12/0/2 & 0{,}009 \\ % E634
\quad GNN (0,04/0,5) & 1{,}299 & \iyi{$-$0{,}101} & [$-$0{,}146; $-$0{,}053] & 12/0/2 & 0{,}005 \\ % E635
\quad Kazanç-MLP & 1{,}156 & \iyi{$-$0{,}245} & [$-$0{,}429; $-$0{,}110] & 12/0/2 & 0{,}001 \\ % E636
\quad Üyelik-MLP & 1{,}207 & \iyi{$-$0{,}193} & [$-$0{,}306; $-$0{,}092] & 11/0/3 & 0{,}004 \\ % E637
\quad Sabit doğrusal & 1{,}162 & \iyi{$-$0{,}239} & [$-$0{,}482; $-$0{,}065] & 10/0/4 & 0{,}020 \\ % E638
\quad Çevrimiçi (ılık) & \kotu{1{,}725} & \kotu{+0{,}325} & [+0{,}150; +0{,}485] & 3/0/11 & 0{,}005 \\ % E639
\addlinespace
\textit{Karışık (14 problem)}: Rastgele & 0{,}645 & \multicolumn{4}{c}{referans} \\ % E640
\quad kNN (0,04/0,75) & 0{,}683 & +0{,}039 & [+0{,}005; +0{,}070] & 3/0/11 & 0{,}042 \\ % E641
\quad Fazla-uzunluk (0,06/0,5) & \iyi{0{,}592} & \iyi{$-$0{,}052} & [$-$0{,}080; $-$0{,}024] & 12/0/2 & 0{,}004 \\ % E642
\quad GNN (0,04/0,5) & 0{,}632 & $-$0{,}013 & [$-$0{,}042; +0{,}016] & 9/0/5 & 0{,}502 \\ % E643
\quad Kazanç-MLP & 0{,}643 & $-$0{,}002 & [$-$0{,}037; +0{,}031] & 8/0/6 & 0{,}952 \\ % E644
\quad Üyelik-MLP & 0{,}637 & $-$0{,}008 & [$-$0{,}031; +0{,}018] & 7/0/7 & 0{,}463 \\ % E645
\quad Sabit doğrusal & 0{,}645 & +0{,}001 & [$-$0{,}034; +0{,}034] & 7/0/7 & 0{,}903 \\ % E646
\quad Çevrimiçi (ılık) & \kotu{0{,}741} & \kotu{+0{,}096} & [+0{,}070; +0{,}124] & 0/0/14 & $1{,}2\times10^{-4}$ \\ % E647
\bottomrule
\end{tabular*}
\end{table}


\paragraph{Vaka incelemesi.}
Şekil~\ref{fig:manset}'te işaretli kümeli h2c02 probleminde GNN hedeflemesi 10~sn sonunda \%1,47 gap'te kalır, fazla-uzunluk \%0,29'a iner.
GNN burada kendi çeyrekler arası aralığındadır; fark rakibin çöküşünden
gelmez. Gerçek eg7146
probleminde gap $2{,}35$'ten $1{,}53$'e, uniform h2u12'de yalnızca $0{,}58$'den $0{,}46$'ya
düşer; yapısal düzensizlik azaldıkça avantaj küçülür.
Karşı-örnek de kümelidir: h2c13'te GNN \%1,53, vekil \%2,16 gap verir.
Kümeli katmanda vekilin GNN'ye karşı farkının güven aralığı sıfırı keser ($-0{,}201$ puan, GA
$[-0{,}439;\,+0{,}014]$); büyük kazanç, büyük varyansla birlikte gelir.
% kaynak: E671, V001-V008

\subsection{Mekanizma}\label{sec:bulgular-mekanizma}

\paragraph{Ucuz ve taze hedef.}
Tablo~\ref{tab:v3-mechanism}, üstünlüğün verimlilikle uyumlu olduğunu gösterir.
Fazla-uzunluk skorunun duvar-saatindeki payı \%0,65'tir; bu pay GNN'de \%8,77, kNN'de
\%2,83'tür.
GNN havuzu seyrek tazelendiği için kesimlerinin \%7,65'i artık turda olmayan kenarlara düşer;
vekilde bu oran \%0,05'tir.
Vekil saniyede toplam 2.288 vuruş yapar (GNN 1.991, kNN 1.640).
Rastgele vuruş daha hızlıdır (3.252), ama kabul oranı \%1,05'te kalır; vekilde bu oran
\%1,41'dir.
Verimlilik tek açıklama değildir: skorlaması GNN'den ucuz olan kNN (model payı \%2,83) ondan iyi
değildir; kNN'nin daha yüksek kabul oranı (\%1,88) vuruş başına yerel arama işini artırır.
% kaynak: E697, E698, E701-E703, E705-E709, E711, E712

\begin{table}[!htbp]
\centering\small\setlength{\tabcolsep}{3pt}
\caption{Mekanizma ve verimlilik, holdout2 (540 koşu/strateji). Model: skor süresinin duvar-saatine oranı; Bayat: turda artık olmayan havuz kenarına düşen kesimler. Renk: sütunun en iyi/en kötüsü.}
\label{tab:v3-mechanism}
\begin{tabular*}{\textwidth}{@{\extracolsep{\fill}}lccccc}
\toprule
Strateji & Vuruş/sn & Kabul [\%] & Model [\%] & Bayat [\%] & Gap [\%] \\
\midrule
Rastgele & 3.252 & 1{,}05 & 0{,}00 & 0{,}00 & 1{,}094 \\
kNN (0,04/0,75) & 1.640 & 1{,}88 & 2{,}83 & 0{,}18 & 1{,}015 \\
Fazla-uzunluk (0,06/0,5) & 2.288 & 1{,}41 & \iyi{0{,}65} & \iyi{0{,}05} & \iyi{0{,}922} \\
GNN (0,04/0,5) & 1.991 & 1{,}67 & \kotu{8{,}77} & \kotu{7{,}65} & 1{,}041 \\
Kazanç-MLP & 1.736 & 1{,}80 & 4{,}86 & 0{,}65 & 1{,}005 \\
Üyelik-MLP & 1.737 & 1{,}86 & 4{,}73 & 0{,}65 & 1{,}011 \\
Sabit doğrusal & 1.751 & 1{,}81 & 4{,}22 & 0{,}40 & 1{,}012 \\
Çevrimiçi (ılık) & 4.357 & 0{,}75 & 4{,}30 & 0{,}07 & \kotu{1{,}303} \\
\midrule
\multicolumn{6}{l}{\textit{Çevrimiçi ağırlık kayması (54 problem, tohum 0): ortalama / medyan / [Q1; Q3]}} \\
\quad Uzunluk: tüm tur kenarları & 0{,}623 & 0{,}609 & \multicolumn{2}{c}{[0{,}602; 0{,}653]} & \\
\quad Uzunluk: ön-bilgi havuzu & 1{,}654 & 1{,}431 & \multicolumn{2}{c}{[1{,}296; 1{,}954]} & \\
\quad Uzunluk: fazla-uz. havuzu & 1{,}614 & 1{,}388 & \multicolumn{2}{c}{[1{,}264; 1{,}918]} & \\
\quad Uzunluk: öğrenilmiş havuz & 0{,}482 & 0{,}487 & \multicolumn{2}{c}{[0{,}445; 0{,}527]} & \\
\quad Örtüşme: ön-bilgi $\cap$ fazla-uz. & 0{,}664 & 0{,}647 & \multicolumn{2}{c}{[0{,}581; 0{,}762]} & \\
\quad Örtüşme: öğrenilmiş $\cap$ fazla-uz. & 0{,}008 & 0{,}001 & \multicolumn{2}{c}{[0{,}000; 0{,}010]} & \\
\quad $\cos(w_0, w_{\text{son}})$ & $-$0{,}518 & $-$0{,}605 & \multicolumn{2}{c}{[$-$0{,}733; $-$0{,}311]} & \\
\bottomrule
\end{tabular*}
\end{table}


\paragraph{Çevrimiçi model ne öğreniyor?}
Holdout2'de güncellemeler başlangıç ağırlıklarını tersine çevirir (kosinüs ortalama $-0{,}52$,
50/54 problemde negatif) ve öğrenilmiş havuz kısa kenarları seçer: göreli uzunluk (kenar
uzunluğunun, uç noktaların 5. en yakın komşu uzaklıkları ortalamasına oranı) öğrenilmiş havuzda
$0{,}48$, ön-bilgi havuzunda $1{,}65$'tir (Şekil~\ref{fig:v3-online-drift}).
Kabul oranı sabit modelin \%1,81'ine karşı \%0,75'tir ve 54 problemin hepsinde düşüktür.
Tablo~\ref{tab:v3-online-reward} bunun nedenini geliştirme setinde ölçer (keşifsel). Öğrenici kendi
amacına ulaşır: vuruş başı ortalama ödülü 12 problemin 12'sinde sabit modelden yüksektir.
Bunu, yerel aramanın aynı maliyete onardığı kısa kenarlara yönelerek yapar: yerel arama öncesindeki
anlık hasar iki modelde aynıyken nötr vuruşların payı \%47'den \%88'e çıkar.
İyileştiren vuruşlar seyrekleşir (\%0,31'den \%0,13'e), kabul edilen bir iyileşmenin ödülü ise
değişmez.
Yani ödül (Denklem~\eqref{eq:reward}) son kaliteyle hizalı değildir. Ödül işareti birim testleriyle
doğrulandığından örüntü bir işaret hatasından kaynaklanmaz.
Geliştirmedeki son gap farkı bu 12 problemde anlamlı değildir ($p=0{,}064$): geliştirme koşuları
mekanizmayı, holdout2 (V1) zararı gösterir.
% kaynak: E599, E722, E726, E730, E732, E736, E737, E938, E941, E943, E944, E947, E948, E959

\begin{table}[!htbp]
\centering\small\setlength{\tabcolsep}{3pt}
\caption{Çevrimiçi ve sabit doğrusal skorun vuruş düzeyi ödülü (geliştirme, 12 VLSI problemi $\times$ 3 tohum, 10~sn; keşifsel). Ç$>$S: çevrimiçi değeri büyük problem sayısı.}
\label{tab:v3-online-reward}
\begin{tabular*}{\textwidth}{@{\extracolsep{\fill}}lcccc}
\toprule
Büyüklük & Çevrimiçi & Sabit & Ç$>$S & Wilcoxon $p$ \\
\midrule
Vuruş başı ödül & $-$2{,}80 & $-$25{,}88 & 12/12 & $4{,}9\times10^{-4}$ \\
Nötr vuruş [\%] & 88{,}1 & 47{,}2 & 12/12 & $4{,}9\times10^{-4}$ \\
Üç kesik kenar geri gelir [\%] & 92{,}4 & 42{,}7 & 12/12 & $4{,}9\times10^{-4}$ \\
İyileştiren vuruş [\%] & 0{,}13 & 0{,}31 & 0/12 & $4{,}9\times10^{-4}$ \\
Kötüleştiren vuruş [\%] & 11{,}8 & 52{,}5 & 0/12 & $4{,}9\times10^{-4}$ \\
Anlık hasar (yerel arama öncesi) & 194{,}8 & 194{,}3 & 8/12 & 0{,}910 \\
Ödül | iyileştiren & 1{,}90 & 1{,}93 & 5/12 & 0{,}791 \\
Havuz kenar uzunluğu / medyan & 1{,}00 & 3{,}35 & 0/12 & $4{,}9\times10^{-4}$ \\
\midrule
Son gap [\%] & 0{,}940 & 0{,}763 & 9/12 & 0{,}064 \\
Vuruş / koşu & 123.078 & 55.406 & -- & -- \\
Kabul / koşu & 145{,}1 & 145{,}0 & -- & -- \\
\bottomrule
\end{tabular*}
\end{table}


\begin{figure}[!htbp]
\centering
\begin{tikzpicture}
\begin{groupplot}[group style={group size=1 by 2, vertical sep=1.9cm},
  width=\textwidth, height=5.0cm, xmode=log, xmin=1450, xmax=10500,
  xtick={1500,2000,3000,4000,6000,8000,10000}, xticklabels={1500,2000,3000,4000,6000,8000,10000},
  xlabel={Instance size $n$ (log scale)}, tick label style={font=\scriptsize}, label style={font=\small},
  title style={font=\small, yshift=-1ex},
  ymajorgrids, grid style={black!10},
  legend style={font=\scriptsize, draw=black!40, fill=white, legend columns=2}]
\nextgroupplot[title={(a) Which edges does the kick pool target?}, ylabel={Relative edge length},
  ymin=0, ymax=3.9, legend style={at={(0.99,0.97)}, anchor=north east}]
\addplot[only marks, mark=square*, mark size=1.4pt, black!75] table[x=n, y=rel_prior] {figures/data/v3_online_drift.dat};
\addlegendentry{Prior weights $w_0$ (static)}
\addplot[only marks, mark=*, mark size=1.4pt, iyiRenk] table[x=n, y=rel_excess] {figures/data/v3_online_drift.dat};
\addlegendentry{Excess proxy}
\addplot[only marks, mark=triangle*, mark size=1.7pt, kotuRenk] table[x=n, y=rel_learned] {figures/data/v3_online_drift.dat};
\addlegendentry{Learned weights $w_{\mathrm{final}}$}
\addplot[only marks, mark=x, mark size=1.4pt, black!40] table[x=n, y=rel_all] {figures/data/v3_online_drift.dat};
\addlegendentry{All tour edges}
\nextgroupplot[title={(b) Kick acceptance rate}, ylabel={Accepted kicks [\%] (log)}, ymode=log, ymin=0.1, ymax=40,
  log ticks with fixed point, legend style={at={(0.01,0.97)}, anchor=north west}]
\addplot[only marks, mark=square*, mark size=1.4pt, black!75] table[x=n, y=acc_linprior] {figures/data/v3_online_drift.dat};
\addlegendentry{Static-linear}
\addplot[only marks, mark=*, mark size=1.4pt, iyiRenk] table[x=n, y=acc_excess] {figures/data/v3_online_drift.dat};
\addlegendentry{Excess}
\addplot[only marks, mark=triangle*, mark size=1.7pt, kotuRenk] table[x=n, y=acc_online] {figures/data/v3_online_drift.dat};
\addlegendentry{Online (warm)}
\addplot[only marks, mark=x, mark size=1.4pt, black!40] table[x=n, y=acc_random] {figures/data/v3_online_drift.dat};
\addlegendentry{Random}
\end{groupplot}
\end{tikzpicture}

\caption{Çevrimiçi (ılık) modelin havuz kayması, holdout2 (tohum 0): (a) havuzların hedeflediği kenarların göreli uzunluğu, (b) kabul edilen vuruş oranı; yatay eksen problem boyutu.}
\label{fig:v3-online-drift}
\end{figure}

\subsection{Duyarlılık: havuz kuralı ve bütçe}\label{sec:bulgular-duyarlilik}

Şekil~\ref{fig:v3-pool-grid}, geliştirme setlerinde önceden yazılmış kuralla seçilen havuz
ayarlarını (Bölüm~\ref{sec:deney-asama3}) ızgara içinde gösterir. Seçilen ayar, kNN, fazla-uzunluk
ve GNN için varsayılandan anlamlı farklı değildir ($p=0{,}82$, $0{,}36$, $0{,}53$).
Izgara $f\ge0{,}04$'te geniş bir plato gösterir. Küçük havuz ($f=0{,}02$) yalnız hedefleme olasılığı
$0{,}75$ ve üzerindeyken geride kalır; olasılık $0{,}9$'da üç sinyal de rastgeleden kötüdür
($+0{,}067$ ile $+0{,}082$ puan arası).
% kaynak: E750-E753, E788

\begin{figure}[!htbp]
\centering
\begin{tikzpicture}
\begin{groupplot}[group style={group size=3 by 1, horizontal sep=0.35cm},
  width=0.385\textwidth, height=5.0cm, xmin=0.01, xmax=0.11, ymin=-0.17, ymax=0.11,
  xtick={0.02,0.04,0.06,0.10}, xticklabels={0.02,0.04,0.06,0.10},
  xlabel={Pool fraction $f$}, tick label style={font=\scriptsize}, label style={font=\small},
  title style={font=\small, yshift=-1ex},
  y tick label style={/pgf/number format/fixed, /pgf/number format/precision=2},
  ymajorgrids, grid style={black!10}]
\nextgroupplot[title={kNN}, ylabel={$\Delta$ gap vs.\ Random [pt]}, legend to name=gridleg, legend columns=4, legend style={font=\scriptsize, draw=black!40}, ]
\addplot[black!45, dashed, forget plot] coordinates {(0.01,0) (0.11,0)};
\addplot[thick, iyiRenk, mark=*, mark size=1.4pt] table[x=frac, y=p050] {figures/data/v3_pool_grid_knn.dat};
\addlegendentry{$p=0.5$}
\addplot[thick, orange!85!black, dashed, mark=triangle*, mark size=1.6pt] table[x=frac, y=p075] {figures/data/v3_pool_grid_knn.dat};
\addlegendentry{$p=0.75$}
\addplot[thick, kotuRenk, densely dotted, mark=square*, mark size=1.3pt] table[x=frac, y=p090] {figures/data/v3_pool_grid_knn.dat};
\addlegendentry{$p=0.9$}
\addplot[only marks, mark=star, mark size=4pt, very thick, black] table[x=frac, y=d] {figures/data/v3_pool_grid_knn_sel.dat};
\addlegendentry{selected rule}
\node[anchor=south east, font=\scriptsize, fill=white, inner sep=1pt, text=black!60] at (axis cs:0.105,0) {Random};
\nextgroupplot[title={Excess}, yticklabels={}, ]
\addplot[black!45, dashed, forget plot] coordinates {(0.01,0) (0.11,0)};
\addplot[thick, iyiRenk, mark=*, mark size=1.4pt] table[x=frac, y=p050] {figures/data/v3_pool_grid_excess.dat};
\addplot[thick, orange!85!black, dashed, mark=triangle*, mark size=1.6pt] table[x=frac, y=p075] {figures/data/v3_pool_grid_excess.dat};
\addplot[thick, kotuRenk, densely dotted, mark=square*, mark size=1.3pt] table[x=frac, y=p090] {figures/data/v3_pool_grid_excess.dat};
\addplot[only marks, mark=star, mark size=4pt, very thick, black] table[x=frac, y=d] {figures/data/v3_pool_grid_excess_sel.dat};
\node[anchor=south east, font=\scriptsize, fill=white, inner sep=1pt, text=black!60] at (axis cs:0.105,0) {Random};
\nextgroupplot[title={GNN}, yticklabels={}, ]
\addplot[black!45, dashed, forget plot] coordinates {(0.01,0) (0.11,0)};
\addplot[thick, iyiRenk, mark=*, mark size=1.4pt] table[x=frac, y=p050] {figures/data/v3_pool_grid_gnn.dat};
\addplot[thick, orange!85!black, dashed, mark=triangle*, mark size=1.6pt] table[x=frac, y=p075] {figures/data/v3_pool_grid_gnn.dat};
\addplot[thick, kotuRenk, densely dotted, mark=square*, mark size=1.3pt] table[x=frac, y=p090] {figures/data/v3_pool_grid_gnn.dat};
\addplot[only marks, mark=star, mark size=4pt, very thick, black] table[x=frac, y=d] {figures/data/v3_pool_grid_gnn_sel.dat};
\node[anchor=south east, font=\scriptsize, fill=white, inner sep=1pt, text=black!60] at (axis cs:0.105,0) {Random};
\end{groupplot}
\node[anchor=north] at ([yshift=-0.9cm]group c2r1.south) {\pgfplotslegendfromname{gridleg}};
\end{tikzpicture}

\caption{Havuz kuralı ızgarası (geliştirme setleri, keşifsel): üç sinyal için havuz oranı $f$ ve hedefleme olasılığına göre rastgeleye fark; yıldız seçilen kuraldır.}
\label{fig:v3-pool-grid}
\end{figure}

Vekilin GNN'ye karşı avantajı erken dönemde büyüktür ve bütçe uzadıkça daralır: fark 0,1~sn'de
$-0{,}44$, 1~sn'de $-0{,}25$, 10~sn'de $-0{,}12$ puandır.
Rastgeleye karşı fark daha kararlıdır (aynı noktalarda $-0{,}17$, $-0{,}28$, $-0{,}17$).
Erken dönemdeki GNN farkının kaynağı model maliyetidir: GNN 0,1~sn'de rastgeleden bile
$0{,}277$ puan kötüdür ve rastgeleden ancak 5~sn'den sonra ayrışır ($-0{,}060$
puan, düzeltmesiz $p=0{,}001$).
Çevrimiçi modelin rastgele vuruşa göre zararı 1~sn'de en büyüktür ($+0{,}369$ puan) ve bütçe
uzadıkça azalarak 10~sn'de $+0{,}209$ puana iner (keşifsel, doğrulayıcı koşunun ara noktaları; V1'deki
$+0{,}291$ puan ise sabit doğrusal modele göre farktır).
% kaynak: E683, E684, E689, E690, E695, E696, E991, E1016, E1026, E1037

Yerel ölçek düzeltmesinin ablasyonu olarak, protokol sonrasında keşifsel bir ham uzunluk kolu (Bölüm~\ref{sec:yontem-sinyal}), rastgele ve
fazla-uzunluk kollarıyla aynı oturumda ve aynı ayarlarla koşuldu. Ham uzunluk rastgeleye
göre gap'i $0{,}189$ puan düşürdü (49/54, $p=1{,}5\times10^{-9}$).
Ölçek düzeltmeli vekille arasında anlamlı bir ortalama fark gözlenmedi ($+0{,}009$ puan, \%95 GA
$[-0{,}032;\,+0{,}053]$, 25'e 28, $p=0{,}99$); güven aralığının $0{,}05$ eşiğine yakınlığı nedeniyle
resmi eşdeğerlik teyidi gelecek ön-kayıtlı bir çalışmaya bırakılmıştır.
Aynı problem ve tohumlarla ayrı oturumda yeniden ölçülen fazla-uzunluk$-$rastgele farkı
$-0{,}197$ puandır (V2: $-0{,}172$).
Ölçek düzeltmesi kümeli problemlerde ortalamada vekil lehine ($0{,}089$ puan; ham kol 14 problemin
6'sında önde), üç ulusal problemde ham uzunluk lehine ($0{,}239$ puan) sonuç verdi; diğer alt
kümelerde farklar $0{,}05$ puanın altındaydı (betimsel).
% kaynak: E600, E965, E972-E979

\subsection{Keşifsel sağlamlık blokları}\label{sec:bulgular-saglamlik}

Tablo~\ref{tab:v4-robustness}, koşulardan önce kaydedilen üç keşifsel bloğun sonuçlarını verir
(Bölüm~\ref{sec:deney-saglamlik}).

\paragraph{Klasik yerel vuruş.}
Hedefli global vuruş, rastgele segment vuruşundan ortalama $1{,}038$ puan düşük gap verir (51/54,
$p_{\mathrm{Holm}}=1{,}3\times10^{-9}$). Fark çarpıktır: en büyük katkı kümeli problemlerden gelir
($-3{,}10$ puan), Hodges--Lehmann kestirimi $-0{,}535$ puandır.
Bu motorda ve bütçede segment geometrisi, rastgele kesimlerle de globalden kötüdür ($+0{,}784$ puan,
7/54); yalnız üç ulusal problemde yerel geometri öndedir.
Segment vuruşu saniyede yaklaşık iki buçuk kat fazla vuruş yapar (5.898'e karşı 2.225), ama kabul
oranı \%0,55'tir (global \%1,41).
Yerel geometride hedeflemenin ölçülebilir etkisi yoktur (fazla-uzunluk $+0{,}037$, ham uzunluk
$-0{,}014$ puan; ikisinde de $p_{\mathrm{Holm}}=1{,}0$). Burada havuzdan yalnız ilk kesim seçilir,
diğer ikisi pencere içinde rastgeledir.
Doğrulama setine veri sızıntısını önlemek amacıyla pencere genişliği ($W=50$) sabit tutulmuştur;
elde edilen sonuç, segment vuruşunun LK tabanlı çözücülerdeki yerleşik performansından
\citep{applegate_2003_clk} bağımsız olarak, 2-opt/Or-opt motorundaki göreli davranışını karakterize eder.
% kaynak: E1047-E1049, E1051, E1057, E1059

\paragraph{Uzun bütçe.}
60~sn'de (5 tohum) vekil hem rastgeleden ($-0{,}071$ puan, 40/54, $p_{\mathrm{Holm}}=7{,}7\times10^{-4}$)
hem GNN'den ($-0{,}078$ puan, 35/54, $p_{\mathrm{Holm}}=0{,}030$) düşük gap verir.
Avantaj bütçe uzadıkça daralır: rastgeleye karşı 10, 30 ve 60~sn'de $-0{,}242$, $-0{,}108$ ve
$-0{,}071$; GNN'ye karşı $-0{,}149$, $-0{,}092$ ve $-0{,}078$ puandır.
GNN'ye karşı farkın \%95 GA'sı 60~sn'de de sıfırın altındadır ($[-0{,}172;\,-0{,}008]$); V3'ün
0,05 puanlık kötü-değil sonucu bu bütçede de geçerlidir.
Mutlak gap'ler de küçüldüğü için göreli kazanç daha yavaş azalır: vekil, rastgelenin gap'ini 10~sn'de
\%21 (\%1,139'dan \%0,897'ye), 60~sn'de \%11 (\%0,649'dan \%0,578'e) düşürür.
% kaynak: E1062-E1065, E1072, E1073, E1077-E1080

\paragraph{Havuz kuralı.}
Geliştirmede seçilen kuralın dört komşusundan üçü holdout2'de de rastgeleden açıkça iyidir
($-0{,}178$ ile $-0{,}221$ puan arası, 49--50/54).
Tek istisna, küçük havuzun her kesimde hedeflendiği birleşimdir ($f=0{,}03$, $p=1$): bu ayar
54 problemin yalnız 11'inde rastgeleden iyidir ve seçilen kuraldan $0{,}313$ puan geridedir.
Seçilen kural, komşularından ya iyidir ya da onlardan ayrışmaz ($+0{,}034$ ile $+0{,}076$ puan;
$f=0{,}03$, $p=0{,}25$ ile fark anlamlı değil).
Bu örüntü geliştirme ızgarasıyla aynıdır (Şekil~\ref{fig:v3-pool-grid}): geniş bir plato ve tek
zararlı bölge olarak küçük havuz ile çok yüksek hedefleme olasılığının birleşimi.
% kaynak: E1086-E1093

\begin{table}[!htbp]
\centering\small\setlength{\tabcolsep}{3pt}
\caption{Keşifsel sağlamlık blokları (holdout2; protokol koşulardan önce kaydedildi). Fark [puan]: negatif değer ilk terim lehine. Önde: ilk terimin önde olduğu problem sayısı. $p_{\mathrm{Holm}}$: blok içinde (b'de her zaman noktasında) Holm düzeltmeli Wilcoxon. Renk: $p_{\mathrm{Holm}}<0{,}05$ ve GA sıfırı dışlar.}
\label{tab:v4-robustness}
\begin{tabular*}{\textwidth}{@{\extracolsep{\fill}}lcccc}
\toprule
Karşılaştırma & Fark & \%95 GA & Önde & $p_{\mathrm{Holm}}$ \\
\midrule
\multicolumn{5}{@{}l}{\textit{(a) Klasik yerel (segment) vuruş --- 10~sn, 10 tohum}} \\
Fazla-uzunluk (global) $-$ rastgele yerel & \iyi{$-$1{,}038} & [$-$1{,}447; $-$0{,}688] & 51/54 & $1{,}3\times10^{-9}$ \\
Fazla-uzunluk (yerel) $-$ rastgele yerel & +0{,}037 & [$-$0{,}028; +0{,}117] & 32/54 & 1{,}000 \\
Ham uzunluk (yerel) $-$ rastgele yerel & $-$0{,}014 & [$-$0{,}079; +0{,}049] & 28/54 & 1{,}000 \\
Fazla-uzunluk: yerel $-$ global & \kotu{+1{,}074} & [+0{,}707; +1{,}505] & 2/54 & $2{,}9\times10^{-9}$ \\
Rastgele: yerel $-$ global & \kotu{+0{,}784} & [+0{,}462; +1{,}155] & 7/54 & $3{,}5\times10^{-6}$ \\
\midrule
\multicolumn{5}{@{}l}{\textit{(b) Uzun bütçe --- 60~sn, 5 tohum (aynı koşuların ara noktaları)}} \\
Fazla-uzunluk $-$ rastgele, 10~sn & \iyi{$-$0{,}242} & [$-$0{,}337; $-$0{,}154] & 50/54 & $1{,}3\times10^{-8}$ \\
Fazla-uzunluk $-$ GNN, 10~sn & \iyi{$-$0{,}149} & [$-$0{,}249; $-$0{,}067] & 41/54 & $7{,}6\times10^{-5}$ \\
Fazla-uzunluk $-$ rastgele, 30~sn & \iyi{$-$0{,}108} & [$-$0{,}165; $-$0{,}051] & 41/54 & $1{,}4\times10^{-5}$ \\
Fazla-uzunluk $-$ GNN, 30~sn & \iyi{$-$0{,}092} & [$-$0{,}185; $-$0{,}022] & 37/54 & 0{,}008 \\
Fazla-uzunluk $-$ rastgele, 60~sn & \iyi{$-$0{,}071} & [$-$0{,}116; $-$0{,}024] & 40/54 & $7{,}7\times10^{-4}$ \\
Fazla-uzunluk $-$ GNN, 60~sn & \iyi{$-$0{,}078} & [$-$0{,}172; $-$0{,}008] & 35/54 & 0{,}030 \\
\midrule
\multicolumn{5}{@{}l}{\textit{(c) Havuz kuralı holdout2'de --- 10~sn, 10 tohum}} \\
$f=0{,}03$, $p=0{,}25$ $-$ rastgele & \iyi{$-$0{,}221} & [$-$0{,}291; $-$0{,}157] & 50/54 & $1{,}3\times10^{-8}$ \\
$f=0{,}03$, $p=1{,}00$ $-$ rastgele & +0{,}059 & [$-$0{,}038; +0{,}146] & 11/54 & 0{,}007 \\
$f=0{,}12$, $p=0{,}25$ $-$ rastgele & \iyi{$-$0{,}178} & [$-$0{,}236; $-$0{,}127] & 50/54 & $8{,}9\times10^{-9}$ \\
$f=0{,}12$, $p=1{,}00$ $-$ rastgele & \iyi{$-$0{,}202} & [$-$0{,}278; $-$0{,}136] & 49/54 & $3{,}2\times10^{-8}$ \\
$f=0{,}03$, $p=0{,}25$ $-$ seçilen kural & +0{,}034 & [$-$0{,}001; +0{,}074] & 21/54 & 0{,}159 \\
$f=0{,}03$, $p=1{,}00$ $-$ seçilen kural & \kotu{+0{,}313} & [+0{,}262; +0{,}360] & 2/54 & $2{,}2\times10^{-8}$ \\
$f=0{,}12$, $p=0{,}25$ $-$ seçilen kural & \kotu{+0{,}076} & [+0{,}036; +0{,}122] & 19/54 & 0{,}005 \\
$f=0{,}12$, $p=1{,}00$ $-$ seçilen kural & \kotu{+0{,}052} & [+0{,}027; +0{,}079] & 14/54 & $4{,}7\times10^{-4}$ \\
\bottomrule
\end{tabular*}
\end{table}


\subsection{Hakem odaklı keşifsel bloklar}\label{sec:bulgular-hakem}

Tablo~\ref{tab:v5-reviewer}, koşulardan önce kaydedilen beş oturumun sonuçlarını verir
(Bölüm~\ref{sec:deney-hakem}).

\paragraph{Yayımlanmış öğrenilmiş skor.}
NeuroLKH'nin ileri geçişi problem başına 0,47~sn ($n=1.689$) ile 2,64~sn ($n=9.882$) arasında
sürdü. Bu süre bütçeden düşüldüğünde vekil ortalamada $0{,}098$ puan öndedir, ancak fark problemler
arasında tutarlı değildir (vekil 20/54 problemde önde, $p_{\mathrm{Holm}}=0{,}58$, medyan fark
NeuroLKH lehine $0{,}033$). Ortalamadaki farkı büyük ölçüde üç büyük ulusal problem yaratır: orada
ağın süresi bütçenin \%20--26'sını alır ve NeuroLKH rastgele vuruşun da gerisine düşer.
Süre bütçeden düşülmediğinde ortalamalar eşittir ($+0{,}009$ puan, \%95 GA $[-0{,}067;\,+0{,}069]$),
fakat NeuroLKH problemlerin çoğunda biraz öndedir (37/54, medyan $0{,}066$ puan,
$p_{\mathrm{Holm}}=0{,}034$). Ortalamayı dengeleyen, vekilin üç ulusal problemdeki üstünlüğüdür
($-0{,}84$ puan). Yani yayımlanmış ve büyük veriyle eğitilmiş bir modelin sinyali, maliyeti
yok sayıldığında, vekilden küçük ama tutarlı bir fark taşır; maliyet sayıldığında bu fark kaybolur.
Ayar eşitliği vekilin diğer rakiplere karşı konumunu değiştirmez. Ayarlı kazanç-MLP varsayılan
kuraldan ayrışmaz ($-0{,}013$ puan, $p_{\mathrm{Holm}}=0{,}58$) ve vekil ayarlı kazanç-MLP'den
($-0{,}070$ puan, 35/54, $p_{\mathrm{Holm}}=0{,}026$), GNN'den ($-0{,}119$ puan, 39/54) ve
rastgeleden ($-0{,}150$ puan, 41/54) düşük gap verir.
% kaynak: E1106-E1113, E1120, E1122

\paragraph{Eşit vuruş sayısı.}
Her kol 20.000 vuruş yaptığında ve model süresi sayılmadığında hedeflemenin değeri açıkça görülür:
bütün sinyaller rastgeleden $0{,}22$--$0{,}30$ puan iyidir ve 54 problemin 49--51'inde öndedir.
Vekilin öğrenilmiş skorlara karşı eşit süredeki üstünlüğü ise bu koşulda kaybolur. Vekil GNN ile
aynı düzeydedir ($-0{,}004$ puan, 25/54); kazanç-MLP ($+0{,}074$ puan) ve NeuroLKH ($+0{,}039$ puan)
vekilden biraz iyidir (32/54 ve 40/54), ancak ikisi de Holm düzeltmesinden sonra anlamlı değildir
($p_{\mathrm{Holm}}=0{,}075$ ve $0{,}051$). Vuruş başına en iyi sinyal eğitimsiz kNN sırasıdır: vekilden
$0{,}088$ puan iyidir (37/54, $p_{\mathrm{Holm}}=0{,}011$).
Vuruş hızları farkın kaynağını gösterir: aynı oturumda vekil saniyede 1.196 vuruş yaparken GNN 524
vuruş yapar ve sürenin \%52'sini modele harcar; kazanç-MLP 919, kNN 867, NeuroLKH 1.073 vuruş/sn'dir.
Eşit süredeki sıralama, büyük ölçüde sinyalin niteliğinden değil maliyetinden doğar.
% kaynak: E1128-E1135, E1137, E1139, E1141, E1143, E1145

\paragraph{Çevrimiçi öğrenmede kredi tanımı.}
Çevrimiçi modelin zararı kredi tanımından gelir. İşaretli ödülle model sabit ön-bilgiden $0{,}335$
puan kötüdür (2/54); ödül $\max(0,r)$ ya da $\mathbb{I}[\text{kabul}]$ olduğunda bu zarar tamamen
kaybolur. İki tanım da işaretli ödülden 54 problemin 54'ünde iyidir ($-0{,}389$ ve $-0{,}338$ puan) ve
sabit ön-bilgiden ayrışmaz ($-0{,}054$ puan, $p_{\mathrm{Holm}}=0{,}13$; $-0{,}003$ puan,
$p_{\mathrm{Holm}}=0{,}62$). İşaretli ödüllü model saniyede 2.745 vuruş yapıp yalnız \%0,8'ini kabul
ettirir; diğer iki tanımda bu değerler yaklaşık 1.560 vuruş/sn ve \%1,4'tür. Bu, modelin ucuz ve
nötr vuruşlara kaydığı mekanizmayla (Bölüm~\ref{sec:bulgular-mekanizma}) tutarlıdır. Doğru krediyle
çevrimiçi öğrenme zarar vermez, ama bu bütçede sabit modele bir şey de katmaz.
% kaynak: E1148-E1153, E1157, E1159, E1161

\paragraph{Tek işçi ve kol sırası.}
Tek işçiyle ve rastgele kol sırasıyla vekil yine rastgeleden ($-0{,}093$ puan, \%95 GA
$[-0{,}182;\,-0{,}012]$) ve GNN'den ($-0{,}057$ puan) düşük gap verir. Yalnız 2 tohum koşulduğu için iki
fark da Holm düzeltmesinden sonra anlamlı değildir. Tek işçide vuruş hızı dört işçili oturumların
yaklaşık 1,7 katıdır (rastgele 3.798'e karşı 2.180 vuruş/sn); dört işçi muhtemelen bellek ve
önbellek paylaşımı yüzünden her kolu birlikte yavaşlatır. Bir partide ilk koşan kol diğerlerinden yaklaşık
\%18 yavaş vurur (Spearman $\rho=0{,}12$). Ana koşuda sıra tohumlar arasında döndürüldüğü için bu
etki büyük ölçüde dengelenmiştir.
% kaynak: E1164, E1165, E1167, E1171, E1127, E1181

\paragraph{LKH-3.}
LKH-3'ün ILS kipinde hedefli vuruş rastgele vuruştan ayrışmaz ($+0{,}001$ puan, \%95 GA
$[-0{,}075;\,+0{,}067]$, 30/54, $p_{\mathrm{Holm}}=0{,}59$). Hedefli kol aynı sürede \%22 daha az deneme
yapar (699'a karşı 897); bunun olası nedeni, havuzun her vuruşta $O(n)$ sürede yeniden kurulmasıdır. Ayrıca ILS kipinin
kendisi LKH'nin varsayılan ayarlarından kötüdür ($+0{,}062$ puan, 10/54). LKH-3, 2-opt/Or-opt
motorundan çok daha düşük gap'lere ulaşır (\%0,26--0,32'ye karşı \%1,07--1,22). Bu düzeyde bir
yerel aramada uzun kenar bilgisi vuruş seçimine ek bir kazanç getirmez.
% kaynak: E1172-E1180

% OTOMATIK URETIM -- evidence/scripts/build_v5.py (elle duzenleme)
\begin{table}[!htbp]
\centering\small\setlength{\tabcolsep}{3pt}
\caption{Hakem odaklı keşifsel bloklar (holdout2; protokol koşulardan önce kaydedildi). Fark [puan]: negatif değer ilk terim lehine. Önde: ilk terimin önde olduğu problem sayısı. $p_{\mathrm{Holm}}$: blok içinde Holm düzeltmeli Wilcoxon. Renk: $p_{\mathrm{Holm}}<0{,}05$ ve GA sıfırı dışlar.}
\label{tab:v5-reviewer}
\begin{tabular*}{\textwidth}{@{\extracolsep{\fill}}lcccc}
\toprule
Karşılaştırma & Fark & \%95 GA & Önde & $p_{\mathrm{Holm}}$ \\
\midrule
\multicolumn{5}{@{}l}{\textit{(a) Eşit süre --- 10~sn, 10 tohum; ayarlı $(f,p)$: kazanç-MLP (0{,}04; 0{,}50), NeuroLKH (0{,}10; 0{,}75)}} \\
Fazla-uzunluk $-$ NeuroLKH (süre sayılır) & $-$0{,}098 & [$-$0{,}210; $-$0{,}008] & 20/54 & 0{,}583 \\
Fazla-uzunluk $-$ NeuroLKH (süre sayılmaz) & +0{,}009 & [$-$0{,}067; +0{,}069] & 17/54 & 0{,}034 \\
NeuroLKH (süre sayılır) $-$ rastgele & $-$0{,}052 & [$-$0{,}117; +0{,}022] & 41/54 & 0{,}002 \\
NeuroLKH (süre sayılmaz) $-$ rastgele & \iyi{$-$0{,}158} & [$-$0{,}214; $-$0{,}107] & 47/54 & $7{,}7\times10^{-7}$ \\
Fazla-uzunluk $-$ kazanç-MLP (ayarlı) & \iyi{$-$0{,}070} & [$-$0{,}128; $-$0{,}018] & 35/54 & 0{,}026 \\
Kazanç-MLP (ayarlı) $-$ kazanç-MLP (varsayılan) & $-$0{,}013 & [$-$0{,}058; +0{,}032] & 31/54 & 0{,}583 \\
Fazla-uzunluk $-$ GNN & \iyi{$-$0{,}119} & [$-$0{,}188; $-$0{,}061] & 39/54 & $8{,}6\times10^{-4}$ \\
Fazla-uzunluk $-$ rastgele & \iyi{$-$0{,}150} & [$-$0{,}218; $-$0{,}089] & 41/54 & $1{,}4\times10^{-5}$ \\
\midrule
\multicolumn{5}{@{}l}{\textit{(b) Eşit vuruş sayısı --- 20.000 vuruş, 5 tohum, model süresi sayılmaz}} \\
Fazla-uzunluk $-$ GNN & $-$0{,}004 & [$-$0{,}076; +0{,}065] & 25/54 & 0{,}633 \\
Fazla-uzunluk $-$ kazanç-MLP (ayarlı) & +0{,}074 & [+0{,}002; +0{,}166] & 22/54 & 0{,}075 \\
Fazla-uzunluk $-$ NeuroLKH & +0{,}039 & [$-$0{,}028; +0{,}106] & 14/54 & 0{,}051 \\
Fazla-uzunluk $-$ kNN sırası & \kotu{+0{,}088} & [+0{,}033; +0{,}156] & 17/54 & 0{,}011 \\
GNN $-$ rastgele & \iyi{$-$0{,}220} & [$-$0{,}306; $-$0{,}141] & 51/54 & $1{,}9\times10^{-7}$ \\
Kazanç-MLP (ayarlı) $-$ rastgele & \iyi{$-$0{,}298} & [$-$0{,}423; $-$0{,}185] & 50/54 & $3{,}8\times10^{-8}$ \\
NeuroLKH $-$ rastgele & \iyi{$-$0{,}262} & [$-$0{,}345; $-$0{,}189] & 49/54 & $1{,}1\times10^{-7}$ \\
Fazla-uzunluk $-$ rastgele & \iyi{$-$0{,}224} & [$-$0{,}294; $-$0{,}157] & 51/54 & $2{,}1\times10^{-8}$ \\
\midrule
\multicolumn{5}{@{}l}{\textit{(c) Çevrimiçi öğrenmede kredi tanımı --- 10~sn, 10 tohum}} \\
Çevrimiçi, $\max(0,r)$ $-$ doğrusal ön-bilgi & $-$0{,}054 & [$-$0{,}111; $-$0{,}004] & 31/54 & 0{,}130 \\
Çevrimiçi, $\mathbb{I}[\text{kabul}]$ $-$ doğrusal ön-bilgi & $-$0{,}003 & [$-$0{,}072; +0{,}062] & 29/54 & 0{,}621 \\
Çevrimiçi, işaretli ödül $-$ doğrusal ön-bilgi & \kotu{+0{,}335} & [+0{,}245; +0{,}432] & 2/54 & $9{,}8\times10^{-10}$ \\
Çevrimiçi, $\max(0,r)$ $-$ çevrimiçi, işaretli ödül & \iyi{$-$0{,}389} & [$-$0{,}494; $-$0{,}292] & 54/54 & $9{,}8\times10^{-10}$ \\
Çevrimiçi, $\mathbb{I}[\text{kabul}]$ $-$ çevrimiçi, işaretli ödül & \iyi{$-$0{,}338} & [$-$0{,}423; $-$0{,}261] & 54/54 & $9{,}8\times10^{-10}$ \\
Doğrusal ön-bilgi $-$ rastgele & $-$0{,}081 & [$-$0{,}146; $-$0{,}023] & 32/54 & 0{,}106 \\
\midrule
\multicolumn{5}{@{}l}{\textit{(d) Tek işçi, rastgele kol sırası --- 10~sn, 2 tohum}} \\
Fazla-uzunluk $-$ rastgele & $-$0{,}093 & [$-$0{,}182; $-$0{,}012] & 29/54 & 0{,}195 \\
Fazla-uzunluk $-$ GNN & $-$0{,}057 & [$-$0{,}128; +0{,}007] & 30/54 & 0{,}355 \\
\midrule
\multicolumn{5}{@{}l}{\textit{(e) LKH-3.0.13 içinde vuruş hedefleme --- 10~sn, 5 tohum}} \\
LKH-3 ILS, hedefli vuruş $-$ LKH-3 ILS, rastgele vuruş & +0{,}001 & [$-$0{,}075; +0{,}067] & 30/54 & 0{,}592 \\
LKH-3 ILS, rastgele vuruş $-$ LKH-3 varsayılan & \kotu{+0{,}062} & [+0{,}018; +0{,}110] & 10/54 & $2{,}9\times10^{-5}$ \\
LKH-3 ILS, hedefli vuruş $-$ LKH-3 varsayılan & +0{,}063 & [$-$0{,}009; +0{,}137] & 15/54 & 0{,}001 \\
\bottomrule
\end{tabular*}
\end{table}


\subsection{Aşama 2: kestirim doğruluğu arama faydası değildir}\label{sec:bulgular-asama2}

Aşama~2'de (38 VLSI problemi, 10~sn) GNN hedeflemesi rastgeleye göre küçük ama anlamlı bir yarar
sağladı ($-0{,}052$ puan, 29/38, $p_{\mathrm{Holm}}=0{,}0004$; Tablo~\ref{tab:v2-primary-tests}).
Eğitimsiz kNN sırası ise GNN'den iyiydi (GNN$-$kNN $+0{,}049$ puan, kNN 29/38 önde,
$p_{\mathrm{Holm}}=0{,}0039$).
Bu sıralama kestirim doğruluğuyla çelişir: GNN, LKH-3 (60~sn) referans turunda olmayan kenarları kNN'den
daha iyi ayırt eder (alıcı işletim karakteristiği eğrisinin altındaki alan, AUC, $0{,}854$'e karşı $0{,}821$; 58 problemin 57'sinde;
Tablo~\ref{tab:v2-mechanism}).
Havuz bileşimi bu çelişkiyle uyumludur. Üyelik için eğitilen skorlar havuza, düzeltildiğinde az
kazandıran incelikli hataları alır; kNN havuzundaki yanlış kenarlar daha pahalıdır (referans tura
göre fazla uzunluk kNN'de $1{,}65$, GNN'de $1{,}30$; kNN 57 problemin 57'sinde önde).
Aynı MLP'yi kazanç hedefiyle eğitmek, kNN karşısındaki kaybın yaklaşık dörtte üçünü kapattı
(keşifsel, deney sonuçlar görüldükten sonra tasarlandı; kalan fark $+0{,}016$ puan, $p=0{,}35$).
Holdout2 bu tabloyu kısmen değiştirir: kazanç-MLP ile üyelik-MLP ayrışmaz ($-0{,}006$ puan,
$p_{\mathrm{Holm}}=1{,}0$), Aşama~2'de rastgeleden iyi olan kNN ($-0{,}101$ puan, 34/38) ise
holdout2'de anlamlılığını yitirir.
Dolayısıyla holdout2 bulguları, hedef uyumu, düşük hesaplama maliyeti ve havuz tazeliğinin tekil
etkenlerden ziyade birlikte çalışan tamamlayıcı mekanizmalar olarak değerlendirilmesi gerektiğini
göstermektedir.
% kaynak: E603, E606, E797, E803, E805, E807, E889, E890, E922, E924, E926, E928, E930

% OTOMATIK URETIM -- evidence/scripts/build_v2.py
\begin{table}[!htbp]
\centering\small\setlength{\tabcolsep}{3pt}
\caption{Birincil sette ön-kayıtlı hipotezler (H1--H4, tek Holm ailesi) ve ikincil karşılaştırmalar. $\Delta$: A$-$B gap farkı [puan], negatif A lehine; TOST sınırı $\pm0{,}10$.}
\label{tab:v2-primary-tests}
\begin{tabular*}{\textwidth}{@{\extracolsep{\fill}}lcccccc}
\toprule
Karşılaştırma (A$-$B) & $\Delta$ & \%95 GA & A/E/B & Wilcoxon $p$ & Holm-$p$ & TOST $p$ \\
\midrule
GNN$-$Rastgele (global) & $-$0{,}052 & [$-$0{,}077; $-$0{,}029] & 29/0/9 & $<10^{-4}$ & 0{,}0004 & 0{,}0002 \\
GNN$-$Rastgele (yerel) & +0{,}019 & [$-$0{,}010; +0{,}047] & 15/0/23 & 0{,}1657 & 0{,}3053 & $<10^{-4}$ \\
GNN$-$kNN (global) & +0{,}049 & [+0{,}020; +0{,}077] & 9/0/29 & 0{,}0013 & 0{,}0039 & 0{,}0007 \\
GNN$-$kNN (yerel) & +0{,}014 & [$-$0{,}008; +0{,}035] & 14/0/24 & 0{,}1526 & 0{,}3053 & $<10^{-4}$ \\
\midrule
kNN$-$Rastgele (global) & $-$0{,}101 & [$-$0{,}126; $-$0{,}077] & 34/0/4 & $<10^{-4}$ & --- & 0{,}5349 \\
kNN$-$Rastgele (yerel) & +0{,}004 & [$-$0{,}025; +0{,}034] & 20/0/18 & 0{,}8522 & --- & $<10^{-4}$ \\
Yerel$-$global (rastgele) & +0{,}362 & [+0{,}321; +0{,}402] & 1/0/37 & $<10^{-4}$ & --- & 1{,}0000 \\
\bottomrule
\end{tabular*}
\end{table}

\begin{table}[!htbp]
\centering\small\setlength{\tabcolsep}{2.5pt}
\caption{LKH-3 referans turunda bulunmayan (``yanlış'') tur kenarlarını ayırt etme, VLSI ortalamaları. Taban: yanlış kenar oranı; Kesinlik: havuzdaki yanlış oranı; $\rho$: GNN--kNN Spearman.}
\label{tab:v2-mechanism}
\begin{tabular*}{\textwidth}{@{\extracolsep{\fill}}lccccccccc}
\toprule
 & Taban & \multicolumn{4}{c}{AUC} & \multicolumn{2}{c}{Kesinlik [\%]} & & Örtüşme \\
\cmidrule(lr){3-6}\cmidrule(lr){7-8}
Tur durumu & [\%] & GNN & kNN & Göreli uz. & Uzunluk & GNN & kNN & $\rho$ & [\%] \\
\midrule
Başlangıç turu (58) & 20{,}5 & 0{,}878 & 0{,}852 & 0{,}816 & 0{,}789 & 85{,}2 & 82{,}1 & 0{,}75 & 65 \\
2 sn ILS turu (58) & 15{,}3 & 0{,}854 & 0{,}821 & 0{,}784 & 0{,}751 & 64{,}6 & 60{,}7 & 0{,}80 & 62 \\
\bottomrule
\end{tabular*}
\end{table}

